\documentclass{article}

\usepackage{amsmath,amssymb}
\usepackage{xurl}
\usepackage[hidelinks]{hyperref}

% Shared commands for notes in docs/paper-gaps/.

\DeclareUnicodeCharacter{2080}{\ifmmode{}_{0}\else\textsubscript{0}\fi}
\DeclareUnicodeCharacter{2081}{\ifmmode{}_{1}\else\textsubscript{1}\fi}
\DeclareUnicodeCharacter{2082}{\ifmmode{}_{2}\else\textsubscript{2}\fi}
\DeclareUnicodeCharacter{2099}{\ifmmode{}_{n}\else\textsubscript{n}\fi}

\newcommand{\Lean}{\textsc{Lean}}
\ExplSyntaxOn
\NewDocumentCommand{\leanid}{m}{
  \ifmmode
    \textsf{\detokenize{#1}}
  \else
    \group_begin:
    \urlstyle{sf}
    \tl_set:Nn \l_tmpa_tl {#1}
    \tl_replace_all:Nnn \l_tmpa_tl {\_} {_}
    \exp_args:NV \path \l_tmpa_tl
    \group_end:
  \fi
}
\ExplSyntaxOff
\providecommand{\tr}{\operatorname{tr}}
\newcommand{\tnleanIssueBase}{https://github.com/LionSR/TNLean/issues/}
\newcommand{\tnleanPrBase}{https://github.com/LionSR/TNLean/pull/}
\newcommand{\tnleanDocBase}{https://sirui-lu.com/TNLean/docs/}
\newcommand{\ghissue}[1]{\href{\tnleanIssueBase#1}{GitHub issue \##1}}
\newcommand{\ghpr}[1]{\href{\tnleanPrBase#1}{PR \##1}}
\newcommand{\doclink}[1]{\href{\tnleanDocBase#1.html}{\path{#1}}}

% Dirac notation and the identity operator, as in blueprint/src/macros/common.tex.
% \providecommand keeps note-local definitions authoritative.
\usepackage{dsfont}
\providecommand{\ket}[1]{|#1\rangle}
\providecommand{\bra}[1]{\langle #1|}
\providecommand{\braket}[2]{\langle #1 | #2 \rangle}
\providecommand{\ketbra}[2]{|#1\rangle\!\langle #2|}
\providecommand{\Id}{\mathds{1}}
\providecommand{\one}{\mathds{1}}
\providecommand{\C}{\mathbb{C}}
\providecommand{\MN}[1]{M_{#1}(\C)}
\providecommand{\MD}{M_D(\C)}

% Verdict marker: \gapnote{<kind>}{<status>}. Kind is the policy
% classification (clarification, local-correction, scope-restriction,
% unfaithful, false-source, open-gap); status is open, wip (elimination
% underway), resolved, or historical. Machine-read by the site index;
% prints a verdict line.
\newcommand{\gapnote}[2]{\par\noindent\textbf{Verdict:} #1 (#2).\par}

\importpaperaux{fbc25-}{2502.20257/main-tnlean}

\title{The Combined Symmetry-Protected Phase of the Controlled-$Z$ Example}
\author{}
\date{2026-11-03}

\begin{document}
\maketitle
\gapnote{scope-restriction}{open}

\section{Source statement}

The controlled-$Z$ example of
Ref.~\cite[lines~3794--3968]{FrancoRubio2025MPUGauging} considers the ring
$U=\prod_j\mathrm{CZ}_{j,j+1}$ of controlled-$Z$ gates and a state generated by
a tensor $A$ with three-dimensional bond. Immediately after the gauging, the
source remarks \cite[line~3970]{FrancoRubio2025MPUGauging} that this state
also has the on-site symmetry $\prod_jZ_j$ and belongs to the nontrivial
symmetry-protected phase of the combined $\mathbb{Z}_2\times\mathbb{Z}_2$ symmetry, as pointed
out in Ref.~\cite{Stephen2024SIMPS}. It then draws the state as a ring of
split-index tensors \cite[lines~3970--4001]{FrancoRubio2025MPUGauging}.

\section{Current formal scope}

Write $S^{ab}=L^aR^b$ for the split-index tensor of the factorization
$A^i=R^iL^i$. The formalization proves the following.
\begin{enumerate}
    \item The amplitude of $A$ on a ring of $N\geq1$ sites is
          $\tr(S^{\sigma_0\sigma_1}S^{\sigma_1\sigma_2}\cdots
              S^{\sigma_{N-1}\sigma_0})$, and $S^{ab}=X^aZ^{ab}$.
    \item The state is invariant under $U$, and under $\prod_jZ_j$, which
          acts on $A$ by conjugation with $\operatorname{diag}(1,-1,-1)$.
    \item On the split-index bond, $Z S^{ab}=(-1)^aS^{ab}Z$ and
          $XS^{ab}=(-1)^{ab}S^{ab}X$. The phase $(-1)^a$ is the factor of
          $\prod_jZ_j$ on the left leg of the split-index tensor, and
          $(-1)^{ab}$ is the controlled-$Z$ phase of its two legs. The
          matrices $Z$ and $X$ anticommute, so they form a projective
          representation of $\mathbb{Z}_2\times\mathbb{Z}_2$ whose class in
          $H^2(\mathbb{Z}_2\times\mathbb{Z}_2,\C^\times)$ is nontrivial.
\end{enumerate}
The last item is
\leanid{ControlledZGauging.simpsTensor_isNontrivialClass}.

The combined symmetry is also realized on the original qubits by an exact
normal matrix product operator family
\leanid{ControlledZGauging.combinedFamily_isNormalRepresentation}.
For $g=(p,q)$ its periodic operator is
$U^q(\prod_jZ_j)^p$, equivalently the diagonal matrix of the products of
the split-leg phases around the ring
\leanid{ControlledZGauging.mpo_combinedFamily}. Its tensors have bond
dimension one for $q=0$ and two for $q=1$. The actual operators are unitary,
and both $\ket{\psi_N(A)}$ and the reference product state
$\ket{0\cdots0}$ are invariant at every positive length. Consequently,
action reductions exist for each of these two states under the same family.
The family's anomaly class vanishes
\leanid{ControlledZGauging.combinedFamily_anomalyClass_eq_zero}; this
three-cocycle statement does not compare the two states' action
two-cocycle classes.

What is not formalized is the identification of this projective class with
the phase: that two states with split-index presentations whose bond
implementations of the combined symmetry carry different classes cannot be
joined by a symmetric gapped path. This is the content of
Ref.~\cite{Stephen2024SIMPS}, which classifies phases protected by symmetries
that act on split-index tensors as on-site symmetries. The formal statement is
therefore the invariant that labels the phase, not the phase statement
itself.

\section{Elimination plan}

The phase statement would follow from a formal version of the classification
of Ref.~\cite{Stephen2024SIMPS} for split-index states, or from the
classification of phases protected by matrix product unitary symmetries
through the action tensors and their L-symbols, applied to the combined
$\mathbb{Z}_2\times\mathbb{Z}_2$ family. The exact normal combined family and
its action reductions for both states are constructed. What remains is a
computation of the class of its L-symbols relative to the invariant product
state $\ket{0\cdots0}$, followed by physical path invariance with the source's
exact hypotheses. The split-index bond calculation alone does not supply
that comparison or the physical path argument.

\bibliographystyle{alpha}
\bibliography{references}

\end{document}
